\section{讨论}

\subsection{训练期机密性与发布期隐私}

Paillier 与差分隐私分别作用于训练过程和发布结果，形成互补的保护机制。Paillier 通过随机化密文与同态聚合保护单站点频率向量及未扰动聚合，站点级差分隐私则通过截断、离散噪声和顺序组合约束最终 Tokenizer 的发布分布。协议匹配无 DP 基线保留公共候选与批量训练路径，有助于区分训练机制变化和随机化对成员风险的影响。实验采用分层验证策略：大规模成员风险与效用评估使用协议等价明文聚合，密码学正确性与开销由独立的真实 2048-bit Paillier 实现测量。

非共谋双服务器模型使噪声明文、单客户端密文与解密能力分离。若服务器共谋，内部计算机密性将受到影响；只要随机化发布机制仍按协议执行，外部 Tokenizer 的站点级差分隐私分布继续由相同机制约束。这一区分说明训练期机密性、发布期隐私与系统信任假设需要在协议分析中分别建模。

\subsection{隐私--效用权衡}

主 SA-DP-BPE 配置采用 $\varepsilon=8$、公共范数分布第 75 百分位截断、$b=32$ 和 $K=1024$，是开发集上成员风险、压缩效率与下游性能综合选择的结果。$\varepsilon=4$ 的 SA-DP-BPE 提供更强随机化，Local-DP-BPE 则以 $\varepsilon=16$ 展示另一种噪声位置和预算配置。不同方法之间的结果共同反映隐私预算、候选域、截断阈值和批量合并对最终性能的综合影响。

本文采用基本顺序组合，使每轮预算分配与总预算之间的关系清晰透明。更紧的隐私会计是改善多轮训练预算利用效率的一条后续方向。

\subsection{扩展方向}

第一，Paillier 明文空间为有限环，而理论双边几何机制定义在整数域上。本文通过 2048-bit 模数和有符号安全区间保持编码、聚合与解码一致；有限模域中的严格机制可采用周期化离散噪声，或将越出安全区间的尾部概率纳入 $(\varepsilon,\delta)$ 会计。

第二，面向服务器共谋和主动参与方，可将门限解密、多方噪声生成以及计数合法性证明纳入同一协议扩展。多语言语料、更大客户端规模、不同预分词和下游任务可合并为统一的外部有效性评估。
